\documentclass[11pt]{article}
\usepackage[margin=1in]{geometry}
\usepackage{graphicx}
\usepackage{setspace}
\usepackage{hyperref}

\title{LIS 470 -- Assignment A4 \\
Defining the Problem}
\author{Team 5 \\
Alejandro De La Torre \\
Angelique Carrillo \\
Helena Huang \\
Yanbin Chen \\
Peixin Yang}
\date{Spring 2026}

\begin{document}

\maketitle

\section*{Introduction}

This document summarizes the insights our team developed after conducting contextual inquiry field research on how students plan and execute grocery shopping. The goal of this assignment is to synthesize our observations into themes, user needs, and a concise point of view that will guide later design work. Rather than jumping immediately to solutions, we focused on identifying the underlying problems students experience when trying to manage groceries within the constraints of student life.

Our team collected detailed field notes during several contextual inquiry sessions with students. These sessions involved observing participants planning grocery trips, comparing prices, organizing lists, and navigating different grocery stores and tools. After collecting the notes, we transformed each observation, quote, and insight into individual notes and collaboratively organized them into clusters using an affinity diagram. This process helped us identify patterns in behavior, motivations, and frustrations related to grocery shopping.

\newpage

\section*{Affinity Diagram}

The affinity diagram below represents the synthesis of our field observations. Each note captures a granular insight from our research and was grouped into clusters based on similarity. Through this collaborative clustering process, several key themes emerged that describe the experiences of students planning and completing grocery shopping tasks.

\begin{center}
\includegraphics[width=\textwidth]{a4_g5_affinity-diagram.pdf}
\end{center}

\newpage

\section*{Themes from the Affinity Diagram}

\subsection*{1. Store Frustrations}
Students frequently encounter friction when interacting with grocery stores themselves. These frustrations include items being out of stock, difficulty locating products, and needing to visit multiple stores to complete a grocery trip. These issues increase the cognitive and time burden of grocery shopping, particularly for students with limited transportation or tight schedules.

\subsection*{2. Behavioral Constraints and Motivations}
Students often rely on routines and self-imposed constraints to manage the complexity of grocery shopping. Budget limits, time pressures, and transportation limitations strongly influence how students plan their shopping behavior. Many students develop structured routines, such as making lists or visiting the same stores regularly, in order to reduce decision fatigue.

\subsection*{3. Demographic Constraints}
Several demographic characteristics shape the grocery shopping experience for students. Many participants are enrolled in SNAP or operate under strict budgets. Transportation access is often limited, meaning that students rely on walking, buses, or nearby stores. These factors significantly influence where and how students shop.

\subsection*{4. Planning the Grocery Trip}
Planning plays a major role in successful grocery shopping. Students frequently plan meals ahead of time, check existing food supplies, and create lists before leaving for the store. This preparation helps them avoid overspending and prevents impulse purchases. Planning also helps students navigate multiple stores more efficiently.

\subsection*{5. Social Context: Roommates}
Roommates introduce both coordination challenges and opportunities. Shared refrigerators, overlapping grocery purchases, and informal food sharing influence how students shop. In many cases, students coordinate purchases through communication with roommates or divide responsibility for certain items.

\subsection*{6. Quality and Brand Preferences}
Students often balance the desire for affordable groceries with preferences for specific brands or food quality. Many participants indicated that they prioritize certain brands they trust, while still trying to remain within budget. This creates trade-offs between cost, perceived quality, and convenience.

\newpage

\section*{User Needs}

Based on the patterns identified in our affinity diagram, we distilled five core user needs. Each need captures a fundamental challenge observed during our contextual inquiry sessions.

\begin{enumerate}

\item Students with limited budgets need a way to understand grocery costs before shopping because overspending can disrupt their weekly food plans.

\item Students need a way to plan grocery trips efficiently because visiting multiple stores and comparing prices is time consuming.

\item Students who rely on SNAP or limited transportation need a way to identify accessible grocery options because not all stores meet their financial or logistical constraints.

\item Students living with roommates need a way to coordinate food purchases because shared spaces often lead to duplicate purchases or food waste.

\item Students need a way to balance food quality with affordability because they want nutritious meals while staying within strict budgets.

\end{enumerate}

\section*{Point of View}

After synthesizing our research insights and user needs, we developed the following point-of-view statement that captures the core problem we observed:

\begin{quote}
Budget-constrained students need simple ways to plan efficient grocery trips without increasing cognitive effort or time.
\end{quote}

This point of view highlights a central tension revealed in our observations: while students want to make thoughtful, cost-effective decisions about groceries, the process of gathering information, comparing stores, and organizing purchases can become mentally taxing. Any design solution should therefore aim to reduce cognitive load while supporting informed decision making.

\end{document}